\section{Zielgruppe}
\label{chap:zielgruppe}

\subsection{Primäre Zielgruppe}
\label{sec:zielgruppe-primaer}

Die geplante App richtet sich an erwachsene Privatpersonen, die einzelne Muttermale oder andere Hautstellen bewusst dokumentieren möchten. Für das MVP konzentrieren wir uns auf deutschsprachige Nutzer mit einem Android-Smartphone. Medizinische Vorkenntnisse sollen für die Bedienung nicht erforderlich sein.

Im Vordergrund stehen zwei Nutzungssituationen:

\begin{itemize}
\item \textbf{Regelmäßige Dokumentation:} Eine Person möchte bestimmte Hautstellen über längere Zeit festhalten. Wichtig sind für sie eine einfache Zuordnung, eine übersichtliche Bildhistorie und Erinnerungen.
\item \textbf{Dokumentation aus einem konkreten Anlass:} Einer Person fällt eine Hautstelle auf, die sie bisher nicht bewusst beobachtet hat. Sie möchte eine Aufnahme speichern und benötigt klare Hinweise dazu, was die App leisten kann und wo ihre Grenzen liegen.
\end{itemize}

Auch ältere Personen gehören ausdrücklich zur Zielgruppe. Die App soll für sie ebenso verständlich und einfach bedienbar sein wie für jüngere Nutzer. Wir legen deshalb keine enge Altersgruppe fest, sondern orientieren uns am Bedürfnis, Hautstellen zu dokumentieren. Für die Bedienung setzen wir lediglich grundlegende Smartphone-Kenntnisse voraus.

\subsection{Anforderungen aus der Zielgruppe}
\label{sec:zielgruppe-anforderungen}

Aus diesen Nutzungssituationen ergeben sich folgende Anforderungen:

\begin{itemize}
\item \textbf{Einfache Sprache:} Begriffe wie \enquote{Läsion} müssen erklärt oder in der Oberfläche durch \enquote{Hautstelle} ersetzt werden.
\item \textbf{Kurzer, geführter Ablauf:} Von der Auswahl einer Hautstelle bis zur gespeicherten Aufnahme sollen möglichst wenige, verständliche Schritte nötig sein.
\item \textbf{Gut erkennbare Bedienelemente:} Gut lesbare Schrift, deutliche Kontraste und ausreichend große Schaltflächen sollen die Bedienung auch bei eingeschränktem Sehvermögen oder motorischen Schwierigkeiten erleichtern.
\item \textbf{Lesbare Ergebnisse:} Risikostufen müssen zusätzlich zur Farbe mit Text erläutert werden. Eine Ampelfarbe allein erklärt weder das Ergebnis noch dessen Unsicherheit.
\item \textbf{Kontrolle über die eigenen Bilder:} Nutzer sollen gespeicherte Aufnahmen leicht wiederfinden und löschen können.
\item \textbf{Verständliche Grenzen:} Wenn eine Aufnahme nicht ausgewertet werden kann, muss die App den Grund nennen, statt trotzdem ein Ergebnis anzuzeigen.
\end{itemize}

\subsection{Abgrenzung}
\label{sec:zielgruppe-abgrenzung}

Hautärzte und medizinische Einrichtungen sind keine primäre Zielgruppe des MVP. Es sind weder eine Praxisverwaltung noch eine direkte ärztliche Beratung innerhalb der App vorgesehen. Angehörige können bei schwer erreichbaren Körperstellen beim Fotografieren helfen; eine eigene Rolle oder ein separates Konto ist dafür zunächst nicht geplant.

Für die KI-Auswertung sieht der Architekturentwurf zusätzliche Einschränkungen nach Alter, Hauttyp und Körperstelle vor. Diese sind als technische Grenzen des geplanten Modells zu verstehen. Ob die konkrete Eingrenzung sachlich ausreicht, muss anhand der verwendeten Daten und der Evaluation geprüft werden. Ein Ausschluss bestimmter Gruppen belegt keine Zuverlässigkeit für die übrigen Nutzer.

\textbf{Die beschriebene Zielgruppe ist die Zielgruppe des Konzepts.} Der Studienprototyp ist für die Demonstration und Erprobung der Abläufe gedacht, nicht für medizinische Entscheidungen im Alltag.
